\documentclass[11pt]{article}
\usepackage[a4paper,margin=1.05in]{geometry}
\usepackage[T1]{fontenc}
\usepackage{palatino}
\usepackage{xcolor}
\usepackage{booktabs}
\usepackage{array}
\usepackage{colortbl}
\usepackage{float}
\usepackage{amsmath,amssymb}
\usepackage{titlesec}
\usepackage{fancyhdr}
\usepackage{enumitem}
\usepackage{tcolorbox}
\usepackage[hidelinks]{hyperref}

\definecolor{bcnavy}{RGB}{11,31,58}
\definecolor{bcgold}{RGB}{176,141,87}
\definecolor{bcgrey}{RGB}{90,98,112}
\definecolor{bcrule}{RGB}{210,214,220}
\definecolor{bcbox}{RGB}{247,241,228}
\definecolor{bcnew}{RGB}{232,240,234}

\titleformat{\section}{\color{bcnavy}\Large\bfseries}{\thesection}{0.7em}{}
\titleformat{\subsection}{\color{bcgold}\normalsize\bfseries}{\thesubsection}{0.7em}{}
\titlespacing*{\section}{0pt}{1.5em}{0.5em}

\pagestyle{fancy}\fancyhf{}
\renewcommand{\headrulewidth}{0pt}
\renewcommand{\footrulewidth}{0.4pt}
\renewcommand{\footrule}{\hbox to\headwidth{\color{bcrule}%
  \leaders\hrule height \footrulewidth\hfill}}
\fancyfoot[L]{\footnotesize\color{bcnavy}\textsc{Bayesian Capital}}
\fancyfoot[C]{\footnotesize\color{bcgrey}Systematic Trading $\cdot$ White Paper}
\fancyfoot[R]{\footnotesize\color{bcgrey}\thepage}

\newtcolorbox{keybox}[1]{colback=bcbox,colframe=bcgold,boxrule=0.6pt,
  title={\color{bcnavy}\bfseries #1},coltitle=bcnavy,fonttitle=\bfseries,
  colbacktitle=bcbox,left=6pt,right=6pt,top=4pt,bottom=4pt}
\newtcolorbox{fixbox}[1]{colback=bcnew,colframe=bcnavy,boxrule=0.6pt,
  title={\color{white}\bfseries #1},coltitle=white,fonttitle=\bfseries,
  colbacktitle=bcnavy,left=6pt,right=6pt,top=4pt,bottom=4pt}

\setlength{\parskip}{0.45em}

\begin{document}

\begin{titlepage}
\centering
\vspace*{3.2cm}
{\color{bcnavy}\fboxsep=22pt\colorbox{bcnavy}{\parbox{3.2cm}{\centering
  {\color{bcgold}\fontsize{34}{34}\selectfont\bfseries BC}\\[6pt]
  {\color{white}\footnotesize\textsc{Bayesian Capital}}}}}\\[2.4cm]
{\color{bcgold}\rule{0.62\textwidth}{0.8pt}}\\[0.5cm]
{\color{bcnavy}\LARGE\bfseries Volatility Harvesting with\\[2mm]
Bayesian and Ornstein--Uhlenbeck Entries}\\[0.35cm]
{\color{bcnavy}\large\bfseries Methodology, Evidence, and the Primacy of Allocation}\\[0.3cm]
{\color{bcgrey}\itshape\small A tranched, premium-capture strategy built on a self-tuning Kalman\\
dip-buyer and a slow mean-reverting hedge sleeve}\\[0.4cm]
{\color{bcgold}\rule{0.62\textwidth}{0.8pt}}\\[2.6cm]
{\color{bcnavy}\bfseries Bayesian Capital}\\[0.15cm]
{\color{bcgrey}Systematic Trading Research}\\[0.15cm]
{\color{bcgrey}Revision 2 --- 4 August 2026}\\[0.5cm]
{\color{bcgrey}\itshape\small White paper --- strictly confidential}
\end{titlepage}

\begin{center}\textbf{Abstract}\end{center}
\begin{quote}\small
We describe the two statistical engines behind our volatility-harvesting book and the evidence we
have gathered on how to deploy them. The book buys short-term weakness and sells a fixed premium
into the bounce; its edge is the repeated capture of that premium as prices oscillate. We drive
entries with two deliberately different estimators. The first is a Bayesian local-linear-trend
Kalman filter that infers a hidden ``fair value'' from noisy closes and bids a volatility-scaled
discount beneath it; every weight in it is derived, not hand-set. The second is an
Ornstein--Uhlenbeck (OU) mean-reverter that fades deviations from a slow long-run mean, and whose
mistakes land on different days --- which is what makes it a genuine hedge rather than a second
copy of the first. Finally, and most importantly, we found that \emph{how we allocate capital}
--- across names, and across the two sleeves --- drives realised out-of-sample performance far
more than the precise parameter values do: re-optimising parameters to recent data reliably
improves the backtest and reliably fails out of sample. The durable edge is structural.

\textbf{Revision 2} corrects one equation. The OU sleeve's uncertainty scale, previously the
standard deviation of the closes about the window mean, measured the trend rather than the
innovation and inflated the entry discount precisely when a name was trending. It is now the
standard deviation of the fitted AR(1) residuals. This is a specification correction rather than
a re-fit, and it is the only change to a signal that has survived our out-of-sample tests.
Section~\ref{sec:sigma} sets out the error, the correction and its evidence.
\end{quote}

\tableofcontents
\newpage

\section{Setup and notation}

Time is indexed by trading periods $t=1,2,\dots,T$. The live book runs on \emph{daily} bars; two
names additionally run a weekly variant (\S\ref{sec:alloc}). The mathematics is period-agnostic,
so we keep a generic period index. For each period we observe a bar and the sub-bars within it:
\begin{itemize}[topsep=2pt,itemsep=1pt]
\item period open $O_t$, close $C_t$, high $\bar H_t$, low $\bar L_t$;
\item intra-period highs $h_{t,d}$ and lows $\ell_{t,d}$, so $\bar H_t=\max_d h_{t,d}$ and
      $\bar L_t=\min_d \ell_{t,d}$.
\end{itemize}
Two derived quantities recur throughout:
\begin{equation}
\text{period range}\quad R_t=\bar H_t-\bar L_t,\qquad
\text{running peak}\quad P_t=\max_{s\le t}\bar H_s .
\end{equation}
The range $R_t$ is the harvest fuel: it scales the model's noise terms, the entry discount, and
--- as we show in \S\ref{sec:alloc} --- the capital each name earns. The shared entry parameters
are the peak cap $c$ (maximum fraction below the all-time high at which we will bid), the sale
premium $\pi$ (the fraction of the previous close we book on exit), and the discount $k$.

\section{Execution and capital mechanics}

Both engines produce a single object each period: a target buy price $B_t$, computed causally
from data up to $t-1$. Execution and capital handling are then shared, so the two engines are
directly comparable and freely blendable.

\subsection{Fills and the premium exit}
The reversion we harvest is fixed in advance as a premium above the entry, $S_t=B_t+\pi_t$ with
$\pi_t=\pi C_{t-1}$. Using intra-period highs and lows,
\begin{align}
\text{buy fills in period }t &\iff \min_d \ell_{t,d}\le B_t,
  & d_t^{\star}&=\min\{d:\ell_{t,d}\le B_t\},\\
\text{sell fills at }S_t &\iff \max_d h_{t,d}\ge S_t
  &&\text{(a same-period sale needs } d\ge d_t^{\star}).
\end{align}
We are filled only on a dip into the discounted bid, and we exit into the reversion at the
premium; a filled-but-unsold position is carried, holding its target $S_t$ until the bounce
arrives or a calendar stop-loss closes it.

\subsection{Compounding and tranching}
The deployed fund $\Phi_t$ buys $\Phi_t/B_t$ shares, so a completed round trip multiplies the
fund by $S/B$. A tranche's terminal return multiple is therefore
\begin{equation}
M=\frac{\Phi_{\text{final}}}{\Phi_{\text{init}}}=\prod_i \frac{S_i}{B_i}.
\end{equation}
Splitting capital $K$ into $n$ equal, one-period-staggered tranches gives total profit
\begin{equation}\label{eq:tranche}
\Pi=\sum_{j=1}^{n}\frac{K}{n}\bigl(M_j-1\bigr)=K\bigl(\bar M-1\bigr),
\qquad \bar M=\frac1n\sum_j M_j,
\end{equation}
so profit depends only on the \emph{average} multiple across start-periods. Idle cash earns
interest at annual rate $r$: a cash period of $\Delta t$ days accrues
$I_t=\Phi_t\,r\,\Delta t/365$, credited monthly. Equation~\eqref{eq:tranche} is the identity that
lets us mix methods and names cleanly: the book's profit is a capital-weighted average of its
tranches' multiples, which is precisely why allocation is a first-order decision.

\section{The Bayesian entry}

We treat fair value not as a formula to be guessed but as a quantity we are \emph{uncertain
about} and hold a belief over --- a belief we update each period as a new close arrives. The
mechanism is: nudge the estimate, weighted by how much we trust the new evidence.

\subsection{Beliefs as Gaussians}
We represent a belief about an unknown number by a Gaussian, captured by the mean $\mu$ and the
variance $V$, with standard deviation $\sigma=\sqrt V$. We use Gaussians because they need only
these two numbers and because combining two Gaussians yields another Gaussian --- so updating
stays simple arithmetic.

\subsection{Bayes' rule in precision form}
For Gaussians, Bayes' rule is cleanest in \emph{precision} --- the inverse variance $1/V$.
Combining evidence adds precision, and the updated mean is the precision-weighted average of the
old mean and the new datum:
\begin{equation}
\frac{1}{V_{\text{post}}}=\frac{1}{V_{\text{prior}}}+\frac{1}{V_{\text{data}}},
\qquad
\mu_{\text{post}}=V_{\text{post}}\left(\frac{\mu_{\text{prior}}}{V_{\text{prior}}}
+\frac{\text{data}}{V_{\text{data}}}\right).
\end{equation}
If we believed fair value was $100\pm4$ and a close prints at $108$, also trusted to $\pm4$, the
posterior is $104$ --- halfway, and more confident. Had the close been unreliable ($\sigma=8$),
the same formula would move us only to $101.6$. This automatic ``trust the more reliable source
more'' is the whole point.

\subsection{The state-space model}
We never observe fair value; we see noisy closes. So we posit a hidden level $m_t$:
\begin{align}
\text{(process)}\quad & m_t=m_{t-1}+\eta_t, & \eta_t &\sim\mathcal N(0,q_t),\\
\text{(observation)}\quad & C_t=m_t+\varepsilon_t, & \varepsilon_t &\sim\mathcal N(0,r_t).
\end{align}

\subsection{The key modelling choice: noise scaled by range}
The single most consequential decision is scaling both noises by the period's range $R_t$,
because a turbulent period both prints a noisier close and gives the underlying value more room
to move:
\begin{equation}
r_t=(\lambda R_t)^2 \quad\text{(observation noise)},\qquad
q_t=(\varphi R_t)^2 \quad\text{(process noise)}.
\end{equation}
$\lambda$ is the fraction of a period's range that is mere noise ($\lambda\approx0.4$--$0.7$);
$\varphi$ is how far fair value can genuinely drift ($\varphi\approx0.15$--$0.20$). Only the
ratio $\varphi/\lambda$ truly matters: it sets the filter's reactivity, and a ratio near
$0.3$--$0.4$ gives an estimate that neither lags badly nor jitters with every print.

\subsection{The period update}
With last period's belief $\mathcal N(\mu_{t-1},V_{t-1})$, each period runs predict-then-update:
\begin{align}
\text{Predict:}\quad & \mu_t^-=\mu_{t-1}, & V_t^-&=V_{t-1}+q_t,\\
\text{Innovation:}\quad & e_t=C_t-\mu_t^-, & K_t&=\frac{V_t^-}{V_t^-+r_t}\in(0,1),\\
\text{Update:}\quad & \mu_t=\mu_t^-+K_te_t, & V_t&=(1-K_t)V_t^-.
\end{align}
The gain is \emph{self-tuning}: in a calm period ($r_t$ small) $K_t\to1$ and we trust the close;
in a wild one $K_t\to0$ and we lean on the prior. This single computed number replaces the
hand-set weights of any fixed formula.

\subsection{From belief to bid: the discount $k$}
For the entry we need the uncertainty of \emph{next} period's price, whose variance has three
independent, additive sources:
\begin{equation}\label{eq:bayessig}
\sigma_t=\sqrt{V_t+q_t+r_t}\qquad\text{(one-step predictive standard deviation)}.
\end{equation}
We bid $k$ of these standard deviations below fair value, floored at the open and the peak cap:
\begin{equation}
B_t^{\text{Bayes}}=\min\bigl(\mu_{t-1}-k\sigma_{t-1},\;O_t,\;(1-c)P_t\bigr).
\end{equation}
Because $\sigma$ widens in turbulent periods, the discount automatically deepens when the model
is less sure --- a genuine credible interval, not a fixed offset. We run $k\approx0.5$--$1.0$.

\subsection{Adding a trend: the local-linear-trend filter}
A pure random walk forecasts ``next period = this period'', so in a sustained rise the fair value
lags and the bid sits too low to fill. We repair this by letting the model also estimate a slope
$b_t$ and project it forward --- mean reversion around a \emph{moving} anchor:
\begin{equation}
m_t=m_{t-1}+b_{t-1}+\eta_t^{(m)},\qquad b_t=b_{t-1}+\eta_t^{(b)},\qquad C_t=m_t+\varepsilon_t,
\end{equation}
with process noises $q^{(m)}=(\varphi_L R_t)^2$ and $q^{(b)}=(\psi R_t)^2$, where
$\varphi_L\approx0.2$ and $\psi\approx0.05$ --- deliberately much smaller, because a trend is
persistent and changes far more slowly than the level itself. Carrying the belief now means
tracking a $2\times2$ variance matrix, and the bid sits below the forecast level pushed one
period ahead by the estimated slope:
\begin{equation}
B_t^{\text{LLT}}=\min\bigl((\hat m_{t-1}+\hat b_{t-1})-k\sigma_{t-1},\;O_t,\;(1-c)P_t\bigr).
\end{equation}
The only change from the random-walk version is the extra term $+\hat b_{t-1}$: one quantity,
learned online, that removes the lag. This is the Bayesian sleeve we run in production.

\section{The Ornstein--Uhlenbeck hedge sleeve}

We run a second entry rule whose explicit purpose is diversification. The mathematics is
period-agnostic, so we keep the same notation.

\subsection{Why a deliberately different signal}
The Bayesian local-linear-trend filter, at heart, \emph{respects the trend}: it carries a slope
and projects it forward, so in a sustained rise it lifts its bid to keep filling. Any second
entry sharing that instinct would be right and wrong on the same days as the first --- its P\&L
correlated, adding return without subtracting much risk. A genuine hedge needs a signal whose
mistakes land on different days. The natural opposite of ``extrapolate the trend'' is ``fade the
deviation'': estimate a slow long-run mean and bet that price is pulled back toward it.

\subsection{The process: a price tethered to a moving mean}
The continuous OU process is the simplest mean-reverting diffusion,
\begin{equation}
dX_t=\theta(\mu-X_t)\,dt+\sigma\,dW_t,
\end{equation}
where $\mu$ is the long-run mean, $\theta>0$ the speed of reversion, and $\sigma\,dW_t$ random
buffeting. Sampled at a fixed step it is an exact AR(1) about the mean:
\begin{equation}\label{eq:ar1}
X_t-\mu=\phi\,(X_{t-1}-\mu)+\varepsilon_t,\qquad \phi=e^{-\theta\Delta t}\in(0,1).
\end{equation}
The persistence $\phi$ carries the whole story: $\phi\to1$ is a random walk, $\phi\to0$ snaps
instantly back. Real prices sit between, and the estimator's job is to measure where.

\subsection{Estimating the ingredients from a rolling window}
Over a trailing window of $W$ trading days ending at $t-1$ we read from the closes
$C_{t-W},\dots,C_{t-1}$:
\begin{align}
\text{window mean}\quad & \bar C_t=\frac1W\sum_{s=t-W}^{t-1}C_s,\\
\text{persistence}\quad & \hat\phi_t=\operatorname{med}\{0,\hat\beta_t,0.99\},\qquad
\hat\beta_t=\frac{\sum (C_s-\bar C)(C_{s-1}-\bar C)}{\sum (C_{s-1}-\bar C)^2}.
\end{align}
Here $\hat\beta_t$ is the ordinary-least-squares slope of each close on the previous close ---
exactly the AR(1) coefficient. The clamp to $[0,0.99]$ is essential, not cosmetic: a raw estimate
below $0$ would assert anti-persistence that makes the rule chase phantom bounces, while a value
at or above $1$ asserts a non-stationary random walk in which the reversion signal is
meaningless. Crucially, these come from a \emph{flat rolling window}, not from the Bayesian
recursion --- a different estimator on a different memory is what makes the OU sleeve a new bet
rather than a re-tuning of the Kalman filter.

\subsection{The scale of the discount: residuals, not levels}\label{sec:sigma}

The third ingredient is the uncertainty that scales the entry discount. Through Revision~1 we
used the standard deviation of the closes about the window mean,
\begin{equation}\label{eq:oldsig}
s_t^{\text{(old)}}=\sqrt{\frac1W\sum_{s=t-W}^{t-1}\bigl(C_s-\bar C_t\bigr)^2 }
\qquad\text{\color{bcgrey}(superseded)} .
\end{equation}

\begin{fixbox}{The error in the superseded scale}
Equation~\eqref{eq:oldsig} measures the dispersion of the price \emph{level} about a flat mean.
On a name with a sustained trend, most of that dispersion \emph{is} the trend. It is not the
scale of the innovation $\varepsilon_t$ in \eqref{eq:ar1}, which is what a discount expressed in
standard deviations is supposed to be measured in.

The consequence runs the wrong way. The buffer $k\,s_t$ inflates precisely when a name is
trending hardest, pushing the bid furthest below the market exactly when the sleeve should be
working --- so fills are lost, and the fills that do occur are crashes rather than reversions.
Measured on the live book, $s_t^{\text{(old)}}$ ran at $8.7\%$ of price on TSM and $14.1\%$ on
VRT, against innovation scales of $2.7\%$ and $3.9\%$: roughly two-thirds of it was trend.
\end{fixbox}

The correct scale is the standard deviation of the residuals of the fitted AR(1) --- what the
reversion model cannot explain. Writing $d_s=C_s-\bar C_t$ for the deviation from the window
mean, the residuals and their scale are
\begin{align}
e_s &= d_s-\hat\phi_t\,d_{s-1}, &
\bar e_t &= \frac{1}{W-1}\sum_{s=t-W+1}^{t-1} e_s, \\
s_t &= \sqrt{\frac{1}{W-1}\sum_{s=t-W+1}^{t-1}\bigl(e_s-\bar e_t\bigr)^2}. \label{eq:newsig}
\end{align}
Both estimators use the same window, the same mean and the same persistence; only the residual
is taken before the dispersion is measured. Note that $\bar e_t$ is linear in the two shifted
window means, so no additional pass over the data is required.

Because $s_t$ in \eqref{eq:newsig} is roughly a third of $s_t^{\text{(old)}}$, the buffer $k$ is
not comparable across the two definitions and must be re-scaled; the fitted values rise from
$\approx0.22$--$0.57$ to $\approx0.20$--$0.75$. This is a change of units, not a re-fit of the
kind \S\ref{sec:noreopt} warns against, and we tested it as such.

\subsection{The forecast and the bid}
The one-step forecast starts at the window mean and adds back only the fraction of today's
deviation that history says tends to persist:
\begin{equation}
f_t=\bar C_t+\hat\phi_t\,(C_{t-1}-\bar C_t),
\end{equation}
and the bid is formed exactly as for the Bayesian entry --- a $k$-sigma discount below the
forecast, floored by the open and the peak cap:
\begin{equation}
B_t^{\text{OU}}=\min\bigl(f_t-k\,s_t,\;O_t,\;(1-c)P_{t-1}\bigr).
\end{equation}
Everything downstream is shared: the fill test, the premium exit, and the capital, commission,
interest, stop-loss and tranching of the earlier sections. Only the location $f_t$ and its
uncertainty $s_t$ are made differently.

\subsection{Why it hedges: the two bids move in opposite directions}
\begin{table}[H]\centering\small
\begin{tabular}{lll}
\toprule
\textbf{After a\dots} & \textbf{LLT-Bayesian bid} & \textbf{OU bid}\\
\midrule
sharp rise & rises (projects the up-slope) & falls (price now above the mean)\\
sharp fall & falls (projects the down-slope) & rises (price now below the mean)\\
quiet drift & tracks the level closely & sits near the slow mean\\
\bottomrule
\end{tabular}
\end{table}
The trend filter chases; the mean-reverter fades. When a rally lifts the Bayesian bid into the
market and it fills near a local top, the same rally has pushed the OU bid below price, so it
steps aside; when a slide drags the Bayesian bid down to catch a falling knife, the OU bid is the
one bidding for the washed-out bounce. This is precisely what drives the two sleeves' return
correlation toward zero.

\subsection{Calibration: the lookback $W$ is the hedge dial}
With the exit, premium and peak cap fixed, the OU entry has two free numbers, the lookback $W$
and the buffer $k$; they trade return against independence, and $W$ is the dominant lever. A
short window ($W\approx12$) is a near-clone: a fast mean hugs the recent price, correlating
$0.76$ with the Bayesian sleeve and \emph{lowering} the combined Sharpe. A long window
($W\approx120$) is a true hedge: the return correlation collapses to $0.09$, effectively
independent, which is the configuration that buys risk reduction.

\subsection{Evidence for the corrected scale}\label{sec:sigmaev}

The correction was tested exactly as any structural change is: the sleeve is scored on verified
fills, marked to market, with the buffer $k$ re-fitted on the training window only and every
other parameter held.

\begin{table}[H]\centering\small
\begin{tabular}{lrrrr}
\toprule
\textbf{OU sleeve alone} & \textbf{eq.\eqref{eq:oldsig}} & \textbf{eq.\eqref{eq:newsig}} &
\textbf{$k$ old} & \textbf{$k$ new}\\
\midrule
RKLB & 171\% & \textbf{199\%} & 0.10 & 0.25\\
MU   & 62\%  & \textbf{83\%}  & 0.20 & 0.75\\
VRT  & 68\%  & 70\%           & 0.20 & 0.40\\
TSM  & 32\%  & 34\%           & 0.05 & 0.20\\
VST  & 54\%  & 54\%           & 0.10 & 0.65\\
\bottomrule
\end{tabular}
\end{table}

\begin{keybox}{Out-of-sample verdict}
Across three expanding folds on each of the five daily names, the corrected scale beat the
superseded one in \textbf{10 of 15 folds}, mean $+14.7$pp. At the book's deployed $75\%$ Bayes
share the effect is smaller, because the sleeve carries a quarter of the capital: the five-name
mean rises from $64\%$ to $71\%$ from the equation alone, and to $77\%$ with the re-scaled
buffers.

\textbf{The qualification.} On the two names run on the weekly model, NVDA and AVGO, the
correction \emph{fails} out of sample --- 2 of 6 folds, mean $-8.2$pp --- despite helping both in
sample. Those are the most persistently trending names in the book, which is where the argument
above predicts the largest gain. Pooled across all seven names the correction stands at
\textbf{12 of 21 folds}, $+8.2$pp. We deploy it on the five names whose own evidence supports it
and record that the mechanism is less clean than the reasoning suggests.
\end{keybox}

\section{Optimising the parameters}

Each name carries ten free parameters,
\begin{equation}
\theta=(\underbrace{\lambda,\varphi_L,\psi,k,\pi,c}_{\text{Bayesian}},\;
\underbrace{W,b_{\text{OU}},\pi_{\text{OU}},c_{\text{OU}}}_{\text{Ornstein--Uhlenbeck}})
\in\Theta\subset\mathbb R^{10},
\end{equation}
confined to a box of economically sensible ranges. We fit them by maximising a backtest
objective, and the choice of \emph{how} to search matters as much as \emph{what} we search for,
because the objective surface is badly behaved.

\subsection{The objectives}
The first maximises annualised return subject to a minimum trade count,
\begin{equation}
\max_{\theta\in\Theta}\;R(\theta)=\Bigl(\frac{\Phi_T(\theta)}{\Phi_0}\Bigr)^{1/Y}-1
\quad\text{s.t.}\quad n_b(\theta)\ge n_{\min},
\end{equation}
where the floor $n_{\min}$ keeps the strategy actively harvesting rather than degenerating to a
handful of large trades. The second tunes the hedge for risk-adjusted return via the annualised
Sharpe ratio of the combined book.

\subsection{Robustness: optimise a ridge, not a spike}
A raw point optimum of a backtest is almost always a fragile spike. We therefore optimise an
\emph{averaged} objective over a small neighbourhood of $\theta$: with $P$ multiplicative
perturbations $\rho^{(p)}$ drawn from $\mathcal U(0.97,1.03)$,
\begin{equation}
\tilde S(\theta)=\tfrac12 S(\theta)+\tfrac12\frac1P\sum_{p=1}^{P}
S\bigl(\Pi_\Theta[\theta\odot\rho^{(p)}]\bigr).
\end{equation}
A knife-edge optimum, whose neighbours perform poorly, is penalised; a plateau is favoured. We
have since made this explicit in reporting as well as in fitting: every planning figure we quote
is the \emph{median of the parameter neighbourhood}, not the value at the fitted point, because
the point is what a ridge fitted to one sample produces and the median is what survives the ridge
moving.

\subsection{Why a derivative-free global search}
$R(\theta)$ and $S(\theta)$ are \textbf{piecewise constant}. A fill is the step indicator
$Z_t=\mathbf 1\{L_t\le B_t(\theta)\}$; as $\theta$ moves, $Z_t$ stays flat until a bid crosses a
day's low, at which point a trade appears or vanishes and cascades through the compounding. Hence
$\nabla_\theta R(\theta)=0$ almost everywhere, with jumps on the fill boundaries. Gradients carry
no information and quasi-Newton methods stall, so we use \textbf{differential evolution}: a
population-based, derivative-free global optimiser, seeded by a Sobol sequence with the incumbent
parameters injected as one member, so the search can never return something worse than the status
quo. A fixed random seed makes every fit reproducible.

\subsection{The limit of all this}
None of the above escapes the deeper fact that a backtest objective is one draw from a noisy
process. The robustness term guards against knife-edge parameters; it does \emph{not} guard
against the sample itself being favourable. That is why the fit is always paired with
out-of-sample validation, and why maximising any in-sample objective fails to beat the incumbent
parameters out of sample. Optimisation locates a good region; validation, not the optimiser,
decides what we trust.

\section{Capital allocation: where we found the durable edge}\label{sec:alloc}

The tranching identity $\Pi=K(\bar M-1)$ makes the book's profit a capital-weighted average of
its tranches' multiples. Two allocation questions therefore sit upstream of every parameter:
\emph{how much capital does each name receive}, and \emph{how is each name's capital split
between the sleeves}. These two decisions govern realised, out-of-sample performance far more
than the precise parameter values do.

\subsection{Across names: equal weighting}
The book is seven names. Five run the daily two-sleeve model; two --- NVDA and AVGO --- run a
weekly single-tranche variant in which the bid is the anchor day's open floored by the peak cap,
and neither sleeve signal is used, having been shown not to earn its place at weekly frequency.

Risk-adjusted weighting, return divided by realised volatility, was tested against
inverse-volatility, Sharpe-proportional and return-proportional schemes. \textbf{Equal weighting
beat all of them by roughly 10pp}, and is what we run: the floor and the cap are both set to
$1/n$, which pins each name to an equal share and leaves the risk-adjusted machinery inert. The
result is not that risk adjustment is wrong in principle but that the return estimates feeding it
are too noisy to improve on the simplest possible rule.

\subsection{Within a name: the Bayes\,/\,OU split}
Each daily name's capital is divided between the two sleeves at a flat \textbf{75:25
Bayes:OU}. Per-name splits were tested and not supported: training windows selected a
Bayes-heavy share for essentially every name, so the extra freedom bought nothing the flat tilt
did not already provide, and the realised out-of-sample gain from the tilt itself is about
$+2$pp against the in-sample curve's $+8$.

The split was re-tested against the corrected OU sleeve of \S\ref{sec:sigma}, on the hypothesis
that the tilt might have been compensating for the mis-specification. It was not. The
full-sample curve does invert --- with the corrected sleeve and re-scaled buffers it prefers
$0\%$ Bayes --- but the walk-forward still selects the deployed share, choosing $100\%$ Bayes as
its median on every training window and beating $75\%$ in only 8 or 9 folds of 15 by about
$1.6$pp. The curve's shape is regime-dependent rather than stable. The split stays where it is.

\subsection{The finding that governs everything: re-optimisation does not generalise}
\label{sec:noreopt}
We tested directly whether parameters should be re-fitted to recent data, across successively
stricter experiments on the live feed.
\begin{enumerate}[topsep=2pt,itemsep=2pt]
\item \textbf{In-sample, it looks great.} Re-fitting each name to maximise return over the recent
      window lifted every name's return, by an average of $+71$ percentage points --- measured on
      the very window we fitted to.
\item \textbf{Fit one year, test the next.} The sign flips: the re-fitted parameters lost out of
      sample in 9 of 10 names, by an average of $-218$pp.
\item \textbf{Walk-forward.} Across rolling folds per name, re-tuned parameters beat the
      incumbent baseline in \textbf{6 of 40 out-of-sample folds}, with zero of ten names showing
      a positive average edge.
\end{enumerate}
On the weekly variant the same discipline gives the same answer: across six separate
re-optimisation experiments, re-fitted parameters have won \textbf{1 fold in 18}.

\begin{keybox}{The edge is structural, not parametric}
Re-optimising parameters to recent data is not merely unhelpful --- it is actively destructive out
of sample, occasionally catastrophically. The strategy's durable value lives in the
\emph{mechanism} (the self-tuning Kalman entry, the deliberately decorrelated OU hedge, the
premium exit, the stop-loss and cash carry) and in \emph{how capital is allocated} across that
mechanism --- not in the third decimal place of any parameter.

The correction in \S\ref{sec:sigma} is the exception that clarifies the rule. It is not a re-fit:
it replaces one estimator with a better-specified one, leaving the fitted parameters where they
were and re-scaling only the buffer whose units changed. It is the one change to a signal that
has survived out-of-sample testing, and it survived because it corrected a specification error
rather than searching for a better number.
\end{keybox}

\begin{keybox}{The one caveat that governs the numbers}
Headline returns are measured over a single, bull-biased history and are optimistic in absolute
terms. We lean on them for \emph{ranking} --- which names and which sleeves are relatively better
--- and deliberately dampen that ranking with caps, floors, the fixed split, and the practice of
quoting the parameter-neighbourhood median rather than the fitted point. The out-of-sample audit
above is the reason we distrust point estimates and widen the safety margin wherever a decision
depends on one.
\end{keybox}

\section{Conclusion}

We drive our volatility-harvesting book with two statistical engines that were built to disagree.
The Bayesian local-linear-trend filter infers a hidden fair value and bids a self-tuning,
volatility-scaled discount beneath it, with no hand-set weights. The Ornstein--Uhlenbeck sleeve
fades deviations from a slow mean and, at a long lookback, misfires on different days --- driving
the return correlation to $0.09$ and, blended in, materially reducing drawdown while raising the
Sharpe ratio. Revision~2 corrects the scale of that sleeve's discount, which had been measuring
trend rather than innovation.

Above both sits the decision we found matters most: the allocation of capital across names and
across the two sleeves. Our out-of-sample testing showed that chasing parameter precision
reliably improves the backtest and reliably fails on unseen data, whereas a disciplined,
equal-weighted allocation over a robust mechanism is what actually survives. The book is
therefore built to be robust first and optimal second --- which, on the evidence we have
gathered, is the only ordering that holds up out of sample.

\vfill
{\color{bcrule}\rule{\textwidth}{0.5pt}}\\[0.2em]
{\footnotesize\itshape\color{bcgrey} Reproducible from the live workbook and a validated Python
replica of the model engine that reproduces its cached results exactly. Same-day round trips are
verified against 1-minute or 5-minute bars; returns are marked to market, with any position open
on the final day valued at that day's close. Figures in-sample unless stated; not a forecast or
an offer.}

\end{document}
